\section{Problem Formulation}
\label{sec:problem_formulation}

\subsection{Domain-Randomized Robust Control}
\label{subsec:general_case}
\textbf{David Example Paragraphs:} We consider the problem of domain-randomized robust control. We denote by $\theta \in \Theta$ a vector of parameters (e.g., mass, dimensions, moments of inertia) that parameterize a family of state dynamics in state $x$: 
\begin{equation}
    x_{t+1} = f(x_t, u_t; \theta),
\end{equation}
where the physical parameters are drawn from a distribution of dynamics instances $\theta \sim \mathcal{D}_\theta$. \\

The objective of domain-randomized robust control is to find a control law, $\pi$, that minimizes expected controller loss with respect to the distribution of physical systems $\mathcal{D}_\theta$. Denoting by $J$ the performance loss, we express the general case as: 
\begin{equation}
    \begin{split}
        \pi^* & := \arg\min_{\pi \in \Pi} \mathbb{E}_{\theta \sim \mathcal{D}_\theta}[J(\pi, \theta)] \\
        \text{s.t.} & \hspace{7mm} x_{t+1} = f(x_t, u_t; \theta)\\ 
        & \hspace{7mm} x_0 \sim \mathcal{D}_{0}
    \end{split}
\end{equation}

\subsection{Linear Domain-Randomized Robust Control}
\label{subsec:linear_case}
% Temporary drafting annotations. Replace the macro body by {} to hide them.
% The report's JMLR class already supplies xcolor and theorem environments.
\providecommand{\foundationsourcenote}[1]{{\normalfont\small\color{gray}%
  \hypersetup{citecolor=gray,linkcolor=gray,urlcolor=gray}%
  [Temporary source note: #1]}}

The course foundation concerns \emph{evaluation of a frozen shared gain}:
learn a quadratic candidate from trajectories, bound its violation risk on
fresh trajectories, and separately check whether it certifies all-state
stability. Controller optimization and subsequent changes to the
randomization law belong to the later stages. The following definitions
and result statements specify the foundation; proofs are deferred to
Section~\ref{sec:theory}. All vector norms are Euclidean;
$\|\cdot\|_2$ and $\|\cdot\|_F$ denote matrix operator and Frobenius norms.

\begin{definition}[Frozen evaluation experiment]
\label{def:foundation_experiment}
Let $\Theta$ be a nonempty finite set of known plants, with
$A(\theta)\in\mathbb R^{n\times n}$ and
$B(\theta)\in\mathbb R^{n\times m}$. For a shared gain
$K\in\mathbb R^{m\times n}$, define
\begin{equation}
\label{eq:foundation_closed_loop}
 x_{t+1}=A(\theta)x_t+B(\theta)u_t=G_K(\theta)x_t,
 \qquad u_t=-Kx_t,\qquad G_K(\theta)=A(\theta)-B(\theta)K.
\end{equation}
Fix $K$, positive integers $H$ (horizon) and $N$ (sample size), and a joint probability
law $\mathcal D_{\mathrm{cert}}$ on
$\Theta\times\mathcal X_0$, where
$\mathcal X_0=\{x:\|x\|_2\le r\}$ and $r>0$, before collecting data.
Draw $z_i=(\theta_i,x_{0,i})\stackrel{\mathrm{i.i.d.}}{\sim}
\mathcal D_{\mathrm{cert}}$ and record $H$ transitions of each trajectory
$\tau_i=\tau_K(z_i)$. Hold $\theta_i$ fixed within its trajectory, so
$x_{t,i}=G_K(\theta_i)^t x_{0,i}$. Set
\begin{equation}
\label{eq:foundation_dynamics_bound}
 M_K=\max_{\theta\in\Theta}\|G_K(\theta)\|_2<\infty.
\end{equation}
Dynamics and observations are exact, deterministic, full-state, and
continuous-action, with no saturation, disturbances, constraints, or early
termination. The independent sample count is $N$, not $NH$.
\foundationsourcenote{Course Chapter~3, Sections~3.1--3.2 (discrete dynamics
and feedback), and Chapter~5, Section~5.6.2 (whole trajectories as samples).
\citet{boffi2021learning}, Section~3.1, sample independent
initial states for one fixed continuous-time system. A joint plant/initial-state
law and a finite discrete horizon are our adaptations; no independence
between the two coordinates of $z_i$ is required.}
\end{definition}

\begin{definition}[Quadratic candidates and normalized decrease]
\label{def:foundation_candidates}
Fix $0<a<b<\infty$, a desired decrease coefficient $\eta>0$, and a fitting
margin $\gamma>0$, all before sampling. Define
\begin{equation}
\label{eq:foundation_candidate_class}
 \mathcal P=\{P=P^\top:aI\preceq P\preceq bI\},
 \qquad V_P(x)=x^\top Px,
 \qquad d=\frac{n(n+1)}2.
\end{equation}
The desired condition is
\begin{equation}
\label{eq:foundation_desired_decrease}
 V_P(G_K(\theta)x)-V_P(x)\le-\eta\|x\|_2^2.
\end{equation}
For a nonzero state and its observed successor define
\begin{equation}
\label{eq:foundation_residual}
 r_P(x,x^+)=\frac{x^{+\top}Px^+-x^\top Px}{\|x\|_2^2}+\eta.
\end{equation}
Then $r_P\le0$ is equivalent to~\eqref{eq:foundation_desired_decrease},
while $r_P\le-\gamma$ requires decrease at least
$(\eta+\gamma)\|x\|_2^2$. At $x=0$, the successor is zero and the desired
condition holds automatically; no normalized residual is evaluated there.
\foundationsourcenote{Course Chapter~3, Theorem~3.10, equations
(3.11)--(3.12), and Example~3.12, equations (3.20)--(3.22), supply the
discrete Lyapunov condition. Boffi Section~3.1.1 uses the unnumbered
continuous-time residual $\nabla V^\top\dot x+\alpha(V)$ in (3.1).
The actual one-step difference, state-weighted decrease, and normalization
are our choices, not an Euler substitution into that residual. The class
specializes Boffi (4.3) with feature map $m(x)=x$ and matrix $Q=P$;
the spectral bounds additionally impose positivity and fix scale.}
\end{definition}

\begin{definition}[Trajectory loss and violation risk]
\label{def:foundation_risk}
For a trajectory $\tau=(\theta,x_0,\ldots,x_H)$ containing a nonzero state
among $x_0,\ldots,x_{H-1}$, set
\begin{equation}
\label{eq:foundation_trajectory_loss}
 s_P(\tau)=\max_{\substack{0\le t<H\\x_t\ne0}}r_P(x_t,x_{t+1}),
 \qquad
 \ell_P(\tau)=\min\{1,\max\{0,1+s_P(\tau)/\gamma\}\}.
\end{equation}
For an all-zero trajectory set $s_P(\tau)=-\infty$ and $\ell_P(\tau)=0$.
Define the population loss, empirical loss, and finite-horizon violation risk
by
\begin{align}
\label{eq:foundation_population_empirical}
 L_K(P)&=\mathbb E_{z\sim\mathcal D_{\mathrm{cert}}}
                  \ell_P(\tau_K(z)),
 &\widehat L_{N,K}(P)&=\frac1N\sum_{i=1}^N\ell_P(\tau_i),\\
\label{eq:foundation_violation_risk}
 e_{K,H}(P)&=\Pr_{z\sim\mathcal D_{\mathrm{cert}}}
                   \{s_P(\tau_K(z))>0\}.
\end{align}
The risk counts a trajectory if \emph{any} of its $H$ transitions violates
the desired condition. A fresh $z$ is independent of the fitting sample.
\foundationsourcenote{Boffi (3.2) asks for a negative fitting margin;
(3.3) defines a worst-time trajectory violation probability. These become
a finite-step maximum and a joint-law probability here. Course Chapter~5,
Definitions~5.1--5.2, supply population and empirical risk; Example~5.14
motivates surrogate losses. Our exact clipped ramp replaces the margin
indicator used at the start of Boffi Section~4. Its positive loss can mean
insufficient slack rather than an actual violation.}
\end{definition}

\begin{definition}[Certificate fitting]
\label{def:foundation_fitting}
Given an optimization tolerance $\xi_{\mathrm{opt}}\ge0$, return a measurable
$\widehat P\in\mathcal P$ satisfying
\begin{equation}
\label{eq:foundation_erm}
 \widehat L_{N,K}(\widehat P)
 \le\inf_{P\in\mathcal P}\widehat L_{N,K}(P)+\xi_{\mathrm{opt}}.
\end{equation}
Only $P$ is fitted; $K$ and the recorded trajectories remain fixed.
Zero empirical loss means all nonzero sampled transitions meet the extra
margin; zero-loss feasibility is not assumed.
\foundationsourcenote{Course Chapter~5, Definition~5.3, is ERM, with
approximate optimization allowed in its discussion. This soft formulation
adapts Boffi's feasibility problem (3.2); it does not assume the feasible
zero-margin-error solution required by their Lemma~4.1. Existence of a
common quadratic certificate is not needed to state the learning problem.}
\end{definition}

\paragraph{Foundation results to establish.}
The next two lemmas specify the regularity and complexity inputs; no
stability assumption on $K$ is needed for these finite-horizon statements.

\begin{lemma}[Loss domination and parameter regularity]
\label{lem:foundation_regular_loss}
For every $P,P'\in\mathcal P$, every nonzero transition $x^+=G_K(\theta)x$,
and every trajectory of Definition~\ref{def:foundation_experiment},
\begin{align}
\label{eq:foundation_residual_lipschitz}
 |r_P(x,x^+)-r_{P'}(x,x^+)|
    &\le(M_K^2+1)\|P-P'\|_F,\\
\label{eq:foundation_loss_lipschitz}
 |\ell_P(\tau)-\ell_{P'}(\tau)|
    &\le L_{\mathrm{par}}\|P-P'\|_F,
 &L_{\mathrm{par}}&=\frac{M_K^2+1}{\gamma}.
\end{align}
Moreover $0\le\ell_P\le1$ and
$\mathbf1\{s_P(\tau)>0\}\le\ell_P(\tau)$, hence
$e_{K,H}(P)\le L_K(P)$.
\foundationsourcenote{Course Chapter~5, Proposition~5.19 and
Section~5.6.2, equation (5.42), use parameter-to-loss Lipschitz control.
Boffi Section~4's joint value/gradient norm and the unnumbered parameter
inequality in arXiv v2 Appendix~B.1 serve the same role. The constant
$(M_K^2+1)/\gamma$ and indicator domination are derived for our normalized
discrete loss, not quoted from Boffi.}
\end{lemma}

\begin{lemma}[Finite covering of the candidate class]
\label{lem:foundation_cover}
Let $\mathcal N_F(\varepsilon,\mathcal P)$ be the smallest cardinality of
a Frobenius $\varepsilon$-net with centers in $\mathcal P$.
For every $\varepsilon>0$,
\begin{equation}
\label{eq:foundation_cover}
 \|P\|_F\le b\sqrt n\quad(P\in\mathcal P),\qquad
 \log\mathcal N_F(\varepsilon,\mathcal P)
 \le d\log\!\left(1+\frac{2b\sqrt n}{\varepsilon}\right).
\end{equation}
\foundationsourcenote{Course Chapter~5, Definition~5.15 and
Propositions~5.18--5.20, define covers and bound Euclidean parameter balls.
Boffi (4.3) counts symmetric matrix parameters; v2 Appendix~B.1 uses the
unnumbered volume bound $\log\mathcal N(u;B_2^k(1))\le k\log(1+2/u)$.
Here the symmetric parameter dimension is $d=n(n+1)/2$, and packing centers
inside $\mathcal P$ preserves admissibility. This does not use the
operator-norm/dual-Sudakov refinement of their Theorem~4.3.}
\end{lemma}

\begin{theorem}[Uniform generalization for the frozen controller]
\label{thm:foundation_generalization}
For the experiment and loss defined above, fix
$\varepsilon>0$ and $\delta\in(0,1)$ before drawing the trajectories, and set
\begin{equation}
\label{eq:foundation_generalization_penalty}
 \Delta_N(\varepsilon,\delta)
 =2L_{\mathrm{par}}\varepsilon+
 \sqrt{\frac{d\log(1+2b\sqrt n/\varepsilon)+\log(2/\delta)}{2N}}.
\end{equation}
With probability at least $1-\delta$ over the $N$ independent trajectories,
simultaneously for all $P\in\mathcal P$,
\begin{equation}
\label{eq:foundation_uniform_bound}
 |L_K(P)-\widehat L_{N,K}(P)|\le\Delta_N(\varepsilon,\delta).
\end{equation}
On the same event, any fitted $\widehat P$ satisfies
\begin{align}
\label{eq:foundation_fitted_risk}
 e_{K,H}(\widehat P)
 &\le\min\{1,\widehat L_{N,K}(\widehat P)+\Delta_N(\varepsilon,\delta)\},\\
\label{eq:foundation_erm_population}
 L_K(\widehat P)
 &\le\inf_{P\in\mathcal P}L_K(P)
        +2\Delta_N(\varepsilon,\delta)+\xi_{\mathrm{opt}},
\end{align}
where the last inequality assumes~\eqref{eq:foundation_erm}.
With $\varepsilon=N^{-1/2}$ and other problem constants fixed,
$\Delta_N=O(\sqrt{(d\log(1+N)+\log(1/\delta))/N})$.
\foundationsourcenote{Course Chapter~5, Lemma~5.6, equation (5.16),
and Theorem~5.8, equation (5.20), supply Hoeffding and a finite union bound;
Proposition~5.19 transfers a parameter cover to losses. Section~5.6.2
supplies trajectory-level sampling. The approach parallels Boffi
Theorem~4.2, (4.2), but replaces Dudley's integral (v2 (B.1)) and their
fast-margin Lemma~4.1 with a two-sided concentration bound on one fixed net.
The two approximation errors give $2L_{\mathrm{par}}\varepsilon$.
This explicit elementary bound, including nonzero fitting loss and
optimization error, is our specialization; it does not reproduce Boffi's
$\widetilde O(d/N)$ or structured-matrix rate.}
\end{theorem}

The conclusion concerns violations over $H$ steps under
$\mathcal D_{\mathrm{cert}}$, not universal or infinite-horizon stability.
Its constants are independent of $H$ for this normalized homogeneous linear
construction; the trajectory loss and risk can still depend on $H$.
Changing $K$, changing the sampling law, or selecting evaluation conditions
adaptively requires fresh justified sampling or a separate transfer argument.

\begin{theorem}[Verified exponential stability]
\label{thm:foundation_verification}
Let $P\in\mathcal P$ (possibly fitted from data), and fix a nonempty
retention set $\Theta_{\mathrm{ret}}\subseteq\Theta$ before inspecting results.
For known matrices define
\begin{equation}
\label{eq:foundation_verified_margin}
 \eta_\theta(P)=\lambda_{\min}
  \!\left(P-G_K(\theta)^\top P G_K(\theta)\right),
 \qquad \eta_* =\min_{\theta\in\Theta_{\mathrm{ret}}}\eta_\theta(P).
\end{equation}
For each retained plant,~\eqref{eq:foundation_desired_decrease} holds for
every $x\in\mathbb R^n$ if and only if $\eta_\theta(P)\ge\eta$.
If $\eta_*>0$, then $0<\eta_*\le b$, every retained fixed-plant closed loop
is globally exponentially stable, and for every initial state and $t\ge1$,
\begin{equation}
\label{eq:foundation_exponential_bound}
 \|x_t\|_2^2\le\frac ba
       \left(1-\frac{\eta_*}{b}\right)^t\|x_0\|_2^2.
\end{equation}
When $\eta_*=b$, each retained loop reaches zero in one step.
Every sublevel set $\mathcal E_c(P)=\{x:x^\top Px\le c\}$, $c\ge0$,
is forward invariant for each retained plant. If also $\eta_*\ge\eta$,
the desired residual never violates on those plants at any horizon; in
particular, $e_{K,H}(P)=0$ if $\mathcal D_{\mathrm{cert}}$ is supported on
$\Theta_{\mathrm{ret}}\times\mathcal X_0$.
\foundationsourcenote{Course Chapter~3, Example~3.12, (3.22), gives
the quadratic drift identity; Theorem~3.10 and Lemma~3.11,
(3.17)--(3.19), give stability and the exponential estimate with
coefficients $a,b,\eta_*$. The eigenvalue test is our exact known-matrix
verification oracle. It is separate from Boffi (3.3)'s statistical error
and does not invoke their Theorem~5.1, (5.1)--(5.3), which instead transfers
small violation measure to a restricted decrease region under additional
incremental-stability and geometric assumptions.}
\end{theorem}

A common quadratic certificate is a restriction beyond individual plant
stability. Positive verification is conditional on the displayed test;
an unattainable fitting margin or a failed test must be reported and does
not establish instability. An initial numerical example may deliberately use
a family known to admit a common certificate. Numerical eigenvalue checks must state
their tolerance; the theorem itself is an exact-arithmetic statement.

\paragraph{Proof tasks and connection to later stages.}
Section~\ref{sec:theory} will establish the two lemmas, apply two-sided
Hoeffding concentration to a deterministic candidate net, extend to the
whole class, and specialize the quadratic Lyapunov comparison for verification.
The theorem statements already specify the constants needed by this route;
sharper complexity terms would require a separate proof and are not assumed.
Certificate fitting and verification leave $K$ unchanged. The subsequent
improvement stage will hold a verified $P$ fixed while changing $K$ and
checking how much decrease margin survives; that update is outside the
foundation's generalization theorem.

For context in those later stages, keep performance distinct from stability.
For fixed $Q,R\succeq0$, let
$J_{H,K}(\theta,x)=\sum_{t=0}^{H-1}(x_t^\top Qx_t+u_t^\top Ru_t)
=x^\top P_H(K,\theta)x$, with $x_0=x$.
Define $S_K^J=\{(\theta,x):\|x\|_2\le r,\ J_{H,K}(\theta,x)\le h\}$.
Coverage is $\mu(S_K^J)$ under a fixed reference probability measure $\mu$;
the training distribution $q$ may change. Net coverage change equals gained
mass minus lost mass. The cost matrix $P_H$ is distinct from a Lyapunov
candidate $P$, and finite-horizon acceptability does not establish stability.

\textbf{Research questions.}
What generalizes when a quadratic certificate is fitted to trajectories of a
frozen controller? Which verified margins survive an evaluation-guided update?
When can targeted training reduce its objective yet lose acceptable coverage?
